\subsection{DEFINITION OF INTEGERS}

DEFINITION 1. A cardinal \(\mathfrak{a}\) is said to be finite if \(\mathfrak{a} \neq \mathfrak{a} + 1\) . A finite cardinal is also called a natural integer (or simply an integer if there is no risk of confusion (*)). A set \(\mathbf{E}\) is said to be finite if Card (E) is a finite cardinal; and Card (E) is then called the number of elements of \(\mathbf{E}\) .

A family (Chapter II, § 3, no. 4) is said to be finite if its index set is finite.

When we say that the number of objects of a certain type is an integer m, we mean that these objects are elements of a finite set whose number of elements is m. A set whose number of elements is m is also called a set of m elements.

PROPOSITION 1. A cardinal a is finite if and only if a + 1 is finite.

For the relations \(\mathfrak{a} = \mathfrak{b}\) and \(\mathfrak{a} + 1 = \mathfrak{b} + 1\) between cardinals \(\mathfrak{a}\) and \(\mathfrak{b}\) are equivalent (§ 3, no. 4, Proposition 8); the relations \(\mathfrak{a} \neq \mathfrak{a} + 1\) and \(\mathfrak{a} + 1 \neq (\mathfrak{a} + 1) + 1\) are therefore equivalent.

It is clear that \(0 \neq 1\) ; hence 0 is an integer. It follows that 1 and 2 are integers. The cardinals \(2 + 1\) and \((2 + 1) + 1\) are integers, denoted by 3 and 4, respectively.

